\section{Introduction}
\label{sec:intro} 
    Federated Learning (FL) \cite{pmlr-v54-mcmahan17a} faces a trilemma \cite{9084352,MAL-083}: client drift from non-IID data \cite{9084352, GUPTA20181}, system heterogeneity, and high communication overhead \cite{MAL-083}. While FedProx \cite{9084352, dai2023tacklingFedProx} targets non-IID data, it still requires full-model training, which is infeasible for resource-constrained devices \cite{10825769}. Moreover, split Learning (SL) \cite{GUPTA20181, vepakomma2019split} solves the computational burden but has a sequential bottleneck \cite{thapa2022splitfed}. Split FL (SFL) \cite{thapa2022splitfed} parallelizes SL but introduces scalability bottlenecks \cite{10980018} and prohibitive communication/privacy overhead from exchanging high-dimensional activations \cite{11091500, zhang2025amperecommunicationefficienthighaccuracysplit, qiu2023evaluating,wang2025towards}. We propose H-SFP (Hierarchical Split-Federated Prototyping) to address these limitations. H-SFP integrates: (i) a hierarchical topology \cite{10980018, 9148862} for scalability; (ii) split execution \cite{GUPTA20181} for low client computation; and (iii) prototype-based communication \cite{tan2022fedproto}, replacing costly activation/gradient exchange \cite{11091500} with lightweight and abstract prototypes. 
    
    This integration yields a crucial synergistic benefit. While FedProto \cite{tan2022fedproto} shows surprising robustness to non-IID data (as our results in \cref{tab:main_accuracy_results} confirm), its flat, centralized architecture does not address server-side scalability. Our work builds on this, positioning H-SFP to not only surpass FedProto's state-of-the-art accuracy but also to simultaneously solve the scalability bottleneck. By performing ``intragroup prototype aggregation" at the edge-server level, H-SFP refines and stabilizes representations, filtering client-level noise before global aggregation. This makes the framework uniquely robust to non-IID data \cite{9084352, GUPTA20181}.

    Our primary contributions are: 
    \begin{itemize} 
        \item We propose H-SFP, a novel three-tier architecture that synergistically integrates split learning, a hierarchical topology, and prototype-based communication to to solve distributed learning bottlenecks holistically. 
        \item We prove hierarchical aggregation acts as a stabilization mechanism, making ours robust to the high statistical heterogeneity challenging standard prototype methods. 
        \item We provide extensive empirical validation of H-SFP's efficiency and scalability against standard FL, HierFL, SFL, HSFL, contrastive-based FL, and FedProto baselines. 
    \end{itemize}